\section{Related Work}\label{sec:related}
\subsection{Document and Diagram Question Answering}
TextVQA introduces questions requiring recognition and contextual use of text in images~\cite{singh2019textvqa}. DocVQA transfers question answering to document pages, where layout and structure determine how textual evidence should be interpreted~\cite{mathew2021docvqa}. ChartQA and InfographicVQA extend the problem to joint interpretation of text, graphics, and quantitative relations~\cite{masry2022chartqa,mathew2022infographic}. These resources establish that reading and reasoning should be assessed together, but their document semantics differ from the conductor and contact relations of an electrical drawing.

AI2D represents diagram elements and their relations for diagram understanding~\cite{kembhavi2016diagram}. GQA uses structured scene representations to generate compositional questions and evaluate consistency~\cite{hudson2019gqa}. These approaches motivate explicit representations and relational checks, although a scene graph does not by itself express electrical continuity under changing contact states. MMMU provides a broader test of multimodal subject knowledge, including engineering~\cite{yue2024mmmu}; a specialized benchmark is needed to resolve failures at individual terminals and connections.

\subsection{Engineering Drawings and Professional Questions}
DrawingVQA is closely related in its use of authentic engineering drawings and professionally meaningful questions~\cite{jung2026drawingvqa}. Its construction-drawing benchmark contains 33 drawings and 92 curated QA pairs, organized by reasoning depth and by engineering and model-capability dimensions. This motivates separating recognition from contextual and engineering judgement. Our setting differs in the target artifact and verification mechanism: electrical continuity and explicit state interventions are represented through bounded CAD-derived graphs. We do not compare benchmark accuracy across these different datasets or equate automatic labels with expert-authored annotations.

DesignQA connects engineering documentation, visual artifacts, and requirement checking~\cite{doris2025designqa}. Its emphasis on applying requirements is relevant to our bounded design-check tasks, where requirements are stated in the question rather than inferred from an unstated standard. Work in IEEE Access on multimodal document segmentation further links technical QA to the organization of local visual elements~\cite{lvov2026segmentation}; our benchmark supplies a fixed local image rather than evaluating document retrieval. MechVQA combines mechanical drawings with recognition, reasoning, and judging tasks and domain adaptation~\cite{kou2026mechvqa}. These studies support a dataset-centered comparative design: the contribution lies in the representation of domain decisions, annotation provenance, and diagnostic evaluation, not merely in the number of image-question pairs.

\subsection{Electrical Schematics and Circuit Reasoning}
Earlier graph-based circuit-sketch recognition addresses the structured interpretation of drawn circuit elements~\cite{liu2020circuitsketch}. CircuitVQA addresses components, values, junctions, and positions in circuit images~\cite{mehta2024circuitvqa}. ElectroVizQA focuses on digital electronics problems~\cite{meshram2024electrovizqa}, whereas EEE-Bench spans multiple electrical and electronic engineering subdomains~\cite{li2025eeebench}. CircuitSense emphasizes deriving symbolic circuit models from visual inputs~\cite{akbari2025circuitsense}; OmniSch focuses on grounding schematic entities and recovering spatially weighted netlist graphs~\cite{lu2026omnisch}. Work on near-sight correction further considers local visual inspection in circuit understanding~\cite{kulkarni2025nsc}. These directions show why a response-level score alone cannot localize whether an error arose from reading, topology extraction, or subsequent reasoning.

\dataset complements these resources with native industrial wiring images, explicit contact assumptions, and brief verifiable explanations. Its tasks concern visible external connections, not unrestricted analog simulation, complete device behavior, or undocumented internal wiring. Table~\ref{tab:related} summarizes the scope and verification emphasis of these resources; it is a descriptive comparison rather than a cross-dataset performance ranking.

\begin{table*}[t]
\centering\small
\caption{Related resources and their reported emphasis. Descriptions concern their task definitions, not results from a common evaluation.}
\label{tab:related}
\begin{tabularx}{\textwidth}{@{}p{23mm}XXXX@{}}\toprule
Resource & Source artifact & Question emphasis & Evidence/verification & Evaluation unit \\\midrule
DrawingVQA~\cite{jung2026drawingvqa} & Construction drawings & Perception to expert judgement & Expert-curated QA / reasoning-depth labels & Drawing / question \\
DesignQA~\cite{doris2025designqa} & Engineering documents and visual artifacts & Requirements and compliance & Document-grounded references & Requirement / question \\
MechVQA~\cite{kou2026mechvqa} & Mechanical drawings & Recognition, reasoning, judging & Semi-automated quality-control pipeline & Drawing / question \\
CircuitVQA~\cite{mehta2024circuitvqa} & Circuit images & Components, values, junctions, positions & Circuit-specific QA annotations & Image / question \\
EEE-Bench~\cite{li2025eeebench} & Electrical/electronics problems & Broad engineering problem solving & Curated problem references & Problem \\
CircuitSense~\cite{akbari2025circuitsense} & Schematics and block diagrams & Perception, analysis, design & Symbolic-equation references & Problem / hierarchy \\
OmniSch~\cite{lu2026omnisch} & PCB schematic diagrams & Grounding and graph construction & Entity / netlist graph annotations & Diagram / structured task \\
\dataset & Native substation DWGs & State-dependent local wiring decisions & CAD graph checks + compact certificates & QA / image; source-grouped split \\
\bottomrule\end{tabularx}
\end{table*}

\subsection{Evidence Supervision and Evaluation}
ScienceQA studies multimodal questions with explanatory supervision~\cite{lu2022scienceqa}, and A-OKVQA provides knowledge-based visual questions with rationales~\cite{schwenk2022aokvqa}. Multimodal-CoT separates rationale generation and answer inference~\cite{zhang2024multimodalcot}; Distilling Step-by-Step uses teacher-generated rationales to supervise smaller language models~\cite{hsieh2023distilling}. These results motivate an investigation of intermediate supervision, but do not establish that additional explanatory tokens improve every schematic task. Our evidence originates from CAD graphs and rules, rather than a documented large-model teacher.

Explanation evaluation requires a separate notion of support. Multimodal explanation research distinguishes textual justification from visual evidence~\cite{park2018explanations}, and ERASER distinguishes rationale alignment from faithfulness~\cite{deyoung2020eraser}. DRAGON evaluates evidence localization over diagrams with the correct answer supplied~\cite{iyengar2026dragon}; our primary task instead requires predicting the answer from an image and question. Generated reasoning may also omit factors that actually influenced a prediction~\cite{turpin2023unfaithful}. We therefore evaluate explicit evidence as a checkable solution certificate, without treating fluent prose as proof of internal reasoning.

Visual instruction tuning and modern open vision--language backbones provide the implementation basis~\cite{liu2023llava,bai2025qwen,li2024onevision}. Position, verbosity, and self-preference effects documented for language-model judges~\cite{zheng2023judge} motivate blinded scoring, structured references, and independent reliability checks.

For generative VQA, LAVE replaces exact string agreement with reference-guided answer assessment~\cite{manas2024lave}. Qwen2.5-32B-Instruct has also been evaluated as an instruction-model judge baseline in prior research~\cite{yu2025risejudge}. We adopt reference-guided assessment with explicit electrical criteria for terminal identities, contact conditions, solution completeness, and admissible alternatives. This is a domain adaptation of the evaluation paradigm; prior judge benchmarks do not establish reliability on electrical drawings, which requires independent human calibration.

Semantic text metrics such as BERTScore compare contextual token representations rather than exact word identity~\cite{zhang2020bertscore}. Electrical QA additionally requires explicit assessment of terminal identities and state-dependent conclusions. G-Eval provides a structured rubric-based evaluation framework~\cite{liu2023geval}, while LLMBar shows that persuasive style can mislead evaluators assessing instruction compliance~\cite{zeng2024llmbar}. Together, these studies motivate domain-specific scoring criteria and validation of evaluator decisions, rather than treating generic textual similarity as a complete measure of electrical correctness.
